\section{Marco de tres etapas: Pre-training, Post-training, Inference-time}

\subsection{Perspectiva unificada: las tres palancas de la formación de capacidad}
El expositor divide el sistema de entrenamiento en tres palancas: el preentrenamiento determina la capacidad de la base y la cobertura de conocimiento, el posentrenamiento determina el estilo de comportamiento y las preferencias de tarea, y el escalamiento en tiempo de inferencia determina la cantidad de cómputo de búsqueda que se puede invertir por solicitud. Las tres son complementarias y presentan diferencias claras en el rendimiento marginal.

\begin{importantbox}{Las tres etapas no son una relación de sustitución, sino una relación de división del trabajo}
\begin{itemize}
    \item \textbf{Pre-training}: aporta una capacidad de representación de amplia cobertura y compresión de conocimiento;
    \item \textbf{Post-training}: orienta la capacidad hacia un comportamiento utilizable (conversación, seguimiento de instrucciones, alineación de preferencias);
    \item \textbf{Inference-time}: invierte cómputo adicional de manera dinámica en el momento del despliegue, para aumentar la tasa de éxito en tareas difíciles.
\end{itemize}
\end{importantbox}

\begin{figure}[H]
\centering
\includegraphics[width=0.9\textwidth]{frame-02-training-stack.jpg}
\caption{Video aproximadamente en 00:11:40: el expositor muestra la pila completa, desde el preentrenamiento hasta el posentrenamiento y el escalamiento en tiempo de prueba}
\end{figure}

\subsection{Descomposición de ingeniería: qué tipo de cuello de botella resuelve cada capa}
\begin{table}[H]
\centering
\renewcommand{\arraystretch}{1.2}
\begin{tabular}{p{2.6cm}p{3.2cm}p{4.0cm}p{3.2cm}}
\toprule
\textbf{Etapa} & \textbf{Objetivo principal} & \textbf{Técnicas típicas} & \textbf{Cuellos de botella comunes} \\
\midrule
Pre-training & Aprender representaciones, comprimir conocimiento, construir la capacidad de base & Proporción de datos, optimizador, programación de la tasa de aprendizaje, detalles de la arquitectura & Límite superior de la calidad de los datos, alto costo de cómputo \\
Post-training & Alinear el comportamiento, reforzar capacidades específicas & SFT, DPO/PPO, RLVR & Conflicto en la distribución de los datos, sesgo del reward \\
Inference-time & Aumentar la tasa de éxito y la robustez en problemas difíciles & Best-of-N, autoconsistencia, iteración de recuperación & Aumento de la latencia y el costo, amplificación de los errores de juicio \\
\bottomrule
\end{tabular}
\caption{Marco de entrenamiento de tres etapas y su división de cuellos de botella}
\end{table}

\begin{knowledgebox}{Por qué el curso enfatiza repetidamente la ``combinación''}
Una sola etapa difícilmente resuelve por sí sola el problema del razonamiento. Por ejemplo, si solo se hace posentrenamiento, se queda limitado por la capacidad de la base; si solo se hace escalamiento en tiempo de inferencia, se queda limitado por la calidad de los candidatos. Por eso, la ruta más realista es hacer una ``combinación acumulable'' según el presupuesto, en lugar de buscar un algoritmo universal.
\end{knowledgebox}

\subsection{Resumen de la sección}
El marco de tres etapas ofrece la perspectiva mínima y completa para analizar los modelos de razonamiento: primero se examina la base, luego la alineación, y por último la asignación de cómputo en el momento del despliegue.


